\documentclass{beamer}
\usepackage[utf8]{inputenc}
\usepackage{dsfont}
\usetheme{Darmstadt}
\usecolortheme{default}

\begin{document}



%---
%This block of code defines the information to appear in the Title page
\title[About Beamer] %optional
{Does Urbanisation Predict Election Outcomes?}

\subtitle{A Bayesian's Perspective}

\author[Christy, Cunha, Murthy] % (optional)
{Eldaleona ~Odole \and Leonor ~Cunha \and Anarghya ~Murthy }

\institute[Technische Universität Dortmund] % (optional)
{
  
  Faculty of Statistics\\
}

\date{\today} % (optional)
% {Mid-Term Presentation}

% \logo{\includegraphics[height=0.5cm]{tu.pdf}}

%End of title page configuration block
%---

%---
%The next block of commands puts the table of contents at the beginning of each section and highlights the current section:
\AtBeginSection[]
{
  \begin{frame}
    \frametitle{Table of Contents}
    \tableofcontents[currentsection]
  \end{frame}
}
%---



%The next statement creates the title page.
\frame{\titlepage}

%---
%This block of code is for the table of contents after
%the title page
\begin{frame}
\frametitle{Table of Contents}
\tableofcontents
\end{frame}
%---

%---SECTION 1: INTRO
\section{Introduction}

\begin{frame}
\frametitle{Introduction}
\begin{itemize}
  \item \textbf{Research Question:}
  How does urbanization of a particular district affect result of an election in the US? 
  \item \textbf{Variable of Interest:}
  Winning party in the  House of Representatives 2022 General Election (binary)
\end{itemize}
\end{frame}



%slide
\section{Dataset Description}
\begin{frame}
\frametitle{Dataset}
We wanted to consider different factors in the analysis, with our primary focus being the urbanization of each House district. These factors included: 
\begin{enumerate}
  \item Demographic Data (US Census Bureau)
  \item Urbanization (FiveThirtyEight)
  \item Regional Information (US Census Bureau)
  \item Election Results (FiveThirtyEight)
\end{enumerate}
We combined different sources in order to create our data set containing 16 unique covariates.
\end{frame}  

%slide
% \begin{frame}
% \frametitle{Frame name}

% \begin{footnotesize}
% \begin{verbatim}
% # Set the maximum number of samples
% N <- 10000        
% # Pmf of x
% x <- 1        
% # Initialize a vector to store the final estimates
% estimates <- numeric(N)   
% # Write function to get estimate of the mean
% # for a given sample size 'n'

% get_est_mean <- function(lambda, n){
%   shelf <- numeric(n)    
%   for(i in 1:n){
%     x <- rpois(m, lambda[i])   #simulate i.i.d.
%                                 #Poisson for given lambda
%     shelf[i] <- mean(x)
%   }
%   return(mean(shelf))
% }

% \end{verbatim}
% \end{footnotesize}

% \end{frame}

\section{Model Setup}

%slide
\begin{frame}
\frametitle{----}

\end{frame}


\section{Results - Pooled Model}

%slide
\begin{frame}
\frametitle{=----}

\end{frame}


\section{Results - Unpooled Model}

%slide
\begin{frame}
\frametitle{=----}

\end{frame}

\section{Results - Hierarchical Model}

%slide
\begin{frame}
\frametitle{=----}

\end{frame}


\section{Results - Varying Intercept Model}

%slide
\begin{frame}
\frametitle{=----}

\end{frame}

\section{Centred vs Non-centred Parameterisation}

%slide
\begin{frame}
\frametitle{STAN Code}

\end{frame}




\end{document}
